% =============================================================================
% APPENDICE: Canali e fonti di ispirazione
% =============================================================================

\chapter*{Canali e fonti di ispirazione}
\addcontentsline{toc}{chapter}{Canali e fonti di ispirazione}

Molte delle idee di questo libro sono maturate lontano dalle biblioteche, davanti a uno schermo, guardando video di divulgazione che hanno acceso domande, collegamenti e curiosità. Ringrazio gli autori e i canali che, nel tempo, sono stati per me una fonte costante di ispirazione:

\begin{itemize}
	\item \textbf{Rick Dufer}, per l'ironia con cui affronta la filosofia e il pensiero critico;
	\item \textbf{Kurzgesagt}, per la capacità di rendere visibili idee scientifiche complesse;
	\item \textbf{151eg}, per lo sguardo curioso su tecnologia e cultura digitale;
	\item \textbf{Ennk}, per gli spunti e i punti di vista che hanno alimentato molte riflessioni;
	\item \textbf{Luciano Floridi}, per il pensiero sulla filosofia dell'informazione e dell'intelligenza artificiale;
	\item \textbf{Geopop}, per la divulgazione scientifica chiara e accessibile;
	\item \textbf{Salvatore Sanfilippo}, per la lucidità con cui racconta la programmazione e l'intelligenza artificiale.
\end{itemize}

Ogni errore o semplificazione presente nelle pagine precedenti resta naturalmente mio.
